\documentclass[12pt]{article}
\usepackage[margin=0.75in]{geometry}
\usepackage{fontspec}
\setmainfont{Latin Modern Roman}
\usepackage{microtype}
\usepackage{fancyhdr}
\usepackage{titlesec}
\usepackage{enumitem}
\usepackage{xcolor}
\usepackage[hidelinks]{hyperref}
\setlength{\parindent}{1.1em}
\setlength{\parskip}{0.4em}
\setlength{\headheight}{15pt}
\titleformat{\section}{\large\bfseries\scshape}{}{0pt}{}
\titleformat{\subsection}{\normalsize\bfseries}{}{0pt}{}
\pagestyle{fancy}
\fancyhf{}
\fancyfoot[C]{\thepage}
\renewcommand{\headrulewidth}{0.4pt}

\fancyhead[L]{Whately: One Word, Two Ideas}
\fancyhead[R]{Argument Lab}
\begin{document}
\begin{center}
{\LARGE\bfseries Argument Lab}\par\vspace{0.25em}
{\large\scshape Fixing the Term}\par\vspace{0.5em}
\rule{0.82\textwidth}{0.4pt}
\end{center}

\section*{Part 1: Spot the Shift}
For each argument below, identify the key term that changes meaning and briefly state the two different ideas it refers to.

\begin{enumerate}[leftmargin=*]
\item “A feather is light. What is light cannot be dark. Therefore a feather cannot be dark.”

\item “The judge said the defendant had a fair trial. But the weather that day was not fair—it rained. Therefore the trial was not fair.”

\item “Only natural substances are healthy. This vitamin is made in a laboratory, so it is not natural. Therefore this vitamin is not healthy.”

\item “We must protect free speech. Therefore no one should be allowed to shout down a speaker on campus.”
\end{enumerate}

\section*{Part 2: Repair the Argument}
Choose one argument from Part 1. Rewrite it so that the key term is used in only one clear sense. Then explain whether the repaired argument now supports its conclusion.

\vspace{4cm}

\section*{Part 3: Create}
Write two short arguments on the same topic (your choice).

\begin{itemize}[leftmargin=*]
\item In the first argument, use a key term in two different senses on purpose.
\item In the second argument, use the same term consistently with a fixed meaning.
\end{itemize}

\vspace{3cm}

\section*{Exit Claim}
In one or two sentences, explain how the habit of fixing terms connects to the four-step analysis you learned in Week 1.

\vspace{2cm}

\end{document}